\section{Load}

\paragraph{Input and output.}
The load stage consumes a root one-file AST plus origin-path information and
returns the root program together with the imported Axon files needed by that
program.

\paragraph{Contract.}
Loading is responsible for file discovery and import-root policy.  It does not
merge modules, rewrite references, prune unreachable definitions, or infer
types.  The output is still a collection of one-file ASTs.

\paragraph{Rationale.}
Separating loading from parsing keeps source-text parsing deterministic and
lets resolution operate over an explicit import set.  This also makes it easier
to test imported builtins and model files without hiding file-system effects in
the parser.

\paragraph{Formal view.}
For a root AST \(A_0\), an origin path \(p\), and an import search policy
\(\mathcal{I}\), loading computes a finite import closure:
\[
  \mathsf{load}(A_0,p,\mathcal{I}) = [A_0,A_1,\ldots,A_n].
\]
The result is ordered only for deterministic processing; it is not yet a
single linked program.

\paragraph{Example.}
If a model file imports
\begin{verbatim}
import NN
import Attention (attention, reshape_heads)
\end{verbatim}
load finds the corresponding builtin Axon files and parses them.  The model AST
still contains import declarations, and the imported ASTs remain separate
files.  Qualified names such as \code{Attention.attention} are not resolved
until the next stage.
